% Body Table 2 as of 2026-08-07 (promoted from app:exhibits after the design
% review; it replaced tables/table_geometry_fingerprint.tex, now deleted).
% The nested delta-R^2 ladder: the one-table version of "magnitude first,
% geometry second". Numbers FROZEN (key_numbers.md section 19.1, n=1034,
% run-level, family fixed effects, quarantine-included). CIs recomputed
% 2026-07-30 by acl_analysis/rq1_stats/04_ladder_ci.py for the ladder's OWN
% blocks (B=2000 cell-level bootstrap); the earlier CIs were borrowed from
% 09-Q1's 2-metric shape-unique share, a different estimand.
% No VERIFY-PENDING markers: every value is a frozen/verified anchor.
% The standardized betas and the cluster-robust t moved to app:exhibits
% (the "residual geometry effect" paragraph) to keep the caption to three
% sentences per the body-exhibit rule.
\begin{table}[t]
\centering
\footnotesize
\setlength{\tabcolsep}{4pt}
\caption{How much of the variation in retention each kind of information
explains, over the $1{,}034$ adapters whose update geometry was measured.
Each row adds one kind of information to the row above it: the size of the
update raises the explained share by $0.395$, three measurements of the
update's geometry raise it by a further $0.017$, and knowing which method
produced the update raises it by $0.006$. Intervals are $95\%$ cell-level
bootstraps, and the size step is larger than the geometry step in all
$2{,}000$ replicates.}
\label{tab:ladder}
\resizebox{\columnwidth}{!}{%
\begin{tabular}{l c c c}
\toprule
Information used & $R^2$ & increase & 95\% CI \\
\midrule
Run family only                        & 0.390 & --       & -- \\
$+$ size of the update ($\dw$)          & 0.785 & $+0.395$ & $[+0.309,+0.483]$ \\
$+$ geometry of the update (3 measures) & 0.802 & $+0.017$ & $[+0.007,+0.032]$ \\
$+$ which method produced it            & 0.808 & $+0.006$ & $[+0.003,+0.018]$ \\
\bottomrule
\end{tabular}%
}
\end{table}
